Wigner 小群已经说明：有质量粒子的内部自由度表现为自旋，而无质量粒子的物理内部自由度表现为螺旋度。散射振幅要使用这些量子数，还需要把它们写成能够直接与 Dirac 链或偏振张量缩并的投影算符。

对有质量电子，螺旋度由协变自旋四矢量沿 $\mathbf p$ 平行或反平行的两种取向给出。对真实光子，它对应二维横向偏振平面上的圆偏振本征基。改变外态偏振只改变 LSZ 外腿的投影，不改变 QED 相互作用顶角本身。


\begin{equation}
 \boxed{
 u_h(p)\bar u_h(p)
 =\frac12(\slashed p+m_{\mathrm e} c)
 \left(1+\gamma^5\slashed s_h\right).}
 \label{eq:electron-helicity-projector-dp9}
\end{equation}

\begin{proposition}[螺旋度旋量的秩一投影]
设三动量非零，外态在正能壳上，螺旋度算符与自旋四矢量采用正文约定。则有以下结论（原式 \eqnrefs{eq:electron-helicity-projector-dp9}）：
\begin{equation*}
 \boxed{
 u_h(p)\bar u_h(p)
 =\frac12(\slashed p+m_{\mathrm e} c)
 \left(1+\gamma^5\slashed s_h\right).}
\end{equation*}
\end{proposition}
\begin{proof}
螺旋度把自旋量子化轴选在动量方向，因此仍属于同一自旋自由度的动量适配基。将正文\eqnrefs{eq:spin-fourvector-r9} 中的静止系自旋方向取为 $\boldsymbol\zeta=\sigma_h\widehat{\mathbf p}$，先得到对应的协变自旋四矢量，再代入固定自旋密度矩阵，就得到正文\eqnrefs{eq:electron-helicity-projector-dp9}。



令 $\sigma_h=\pm1$ 表示两倍螺旋度符号，即电子螺旋度本征值为
\begin{equation*}
 h=\frac{\sigma_h}{2}\hbar.
\end{equation*}
对 $\mathbf p\neq0$ 定义 $\widehat{\mathbf p}=\mathbf p/|\mathbf p|$。把静止系自旋方向取为 $\boldsymbol\zeta=\sigma_h\widehat{\mathbf p}$，代入\eqnrefs{eq:spin-fourvector-r9} 的推动公式 \eqnrefs{eq:spin-fourvector-r9}。因为推动与自旋方向共线，得到
\begin{equation*}
 {
 s_h^\mu
 =\sigma_h\left(
 \frac{|\mathbf p|}{m_{\mathrm e} c},
 \frac{E_{\mathbf p}}{m_{\mathrm e} c^2}\widehat{\mathbf p}
 \right).}
 \end{equation*}
逐项检查：
\begin{align*}
 p\!\cdot s_h
 &=\frac{E_{\mathbf p}}{c}\frac{\sigma_h|\mathbf p|}{m_{\mathrm e} c}
 -\mathbf p\cdot
 \frac{\sigma_hE_{\mathbf p}}{m_{\mathrm e} c^2}\widehat{\mathbf p}
 =0,\\
 s_h^2
 &=\frac{|\mathbf p|^2}{m_{\mathrm e}^2c^2}
 -\frac{E_{\mathbf p}^2}{m_{\mathrm e}^2c^4}
 =-1,
\end{align*}
其中最后一步用了 $E_{\mathbf p}^2=m_{\mathrm e}^2c^4+c^2|\mathbf p|^2$。所以它确实是 \eqnrefs{eq:polarized-spinor-projector-r9} 所要求的单位类空自旋四矢量。

直接代入偏振密度矩阵，得到螺旋度可分辨的外部电子投影算符
\begin{equation*}
 {
 u_h(p)\bar u_h(p)
 =\frac12(\slashed p+m_{\mathrm e} c)
 \left(1+\gamma^5\slashed s_h\right).}
 \end{equation*}
把 $\sigma_h=+1$ 与 $-1$ 相加，$\slashed s_h$ 项相消，恢复
$\sum_hu_h\bar u_h=\slashed p+m_{\mathrm e} c$。因此未偏振自旋求和就是两个螺旋度投影算符的完备和。

在 $\mathbf p\to0$ 时 $\widehat{\mathbf p}$ 无定义，这不是公式缺陷：有质量粒子的静止系没有优选传播方向，螺旋度本来就不是静止系量子数。此时应回到任意静止系自旋轴的 $SU(2)$ 基。
\end{proof}


对极端相对论过程，有时还需要区分螺旋度与手征性。手征投影算符已在 \eqnrefs{eq:chiral-projectors} 定义为
\begin{equation*}
 P_{R,L}=\frac12(1\pm\gamma^5).
\end{equation*}
对 $m_{\mathrm e}\neq0$，$P_{R,L}$ 不等于螺旋度投影算符；只有在严格无质量 Dirac 方程或受控 $m_{\mathrm e} c^2/E\to0$ 极限中，两者才沿 positive-能量粒子分支对齐。

\begin{equation}
 \boxed{
 \mathsf P_\lambda^{\mu\nu}(k;n)
 \equiv
 \varepsilon_\lambda^\mu\varepsilon_\lambda^{\nu *}
 =\frac12\left[
 \mathsf P_\perp^{\mu\nu}(k;n)
 -\ii\lambda\mathsf C^{\mu\nu}(k;n)
 \right].}
 \label{eq:photon-helicity-projector-covariant-dp9}
\end{equation}

\begin{proposition}[横向二维空间的螺旋度投影]
设光子四动量为非零零矢量，辅助矢量与其内积非零；极化空间模去纵向规范方向。则有以下结论（原式 \eqnrefs{eq:photon-helicity-projector-covariant-dp9}）：
\begin{equation*}
 \boxed{
 \mathsf P_\lambda^{\mu\nu}(k;n)
 \equiv
 \varepsilon_\lambda^\mu\varepsilon_\lambda^{\nu *}
 =\frac12\left[
 \mathsf P_\perp^{\mu\nu}(k;n)
 -\ii\lambda\mathsf C^{\mu\nu}(k;n)
 \right].}
\end{equation*}
\end{proposition}
\begin{proof}
取零质量光子动量 $k^\mu$，并选任意满足 $k\!\cdot n\neq0$ 的辅助四矢量 $n^\mu$。由\eqnrefs{eq:physical-photon-polarization-projector-dp42} 得到物理横向偏振求和

\begin{equation*}
 \mathsf P_\perp^{\mu\nu}(k;n)
 \equiv
 -\eta^{\mu\nu}
 +\frac{k^\mu n^\nu+n^\mu k^\nu}{k\!\cdot n}
 -\frac{n^2k^\mu k^\nu}{(k\!\cdot n)^2}.
 \tag{\ref{eq:physical-photon-polarization-projector-dp42}}
\end{equation*}
在该二维物理子空间选两个实正交偏振矢量 $e_1^\mu,e_2^\mu$，使
\begin{equation*}
 e_a\!\cdot k=e_a\!\cdot n=0,
 \qquad
 e_a\!\cdot e_b=-\delta_{ab},
 \qquad
 \mathsf P_\perp^{\mu\nu}=e_1^\mu e_1^\nu+e_2^\mu e_2^\nu.
\end{equation*}
定义圆偏振（螺旋度）基
\begin{equation*}
 \varepsilon_\lambda^\mu(k;n)
 \equiv\frac1{\sqrt2}
 \left(e_1^\mu+\ii\lambda e_2^\mu\right),
 \qquad \lambda=\pm1.
 \end{equation*}
于是
\begin{align*}
 \varepsilon_\lambda^\mu\varepsilon_\lambda^{\nu *}
 &=\frac12\Bigl[
 e_1^\mu e_1^\nu+e_2^\mu e_2^\nu
 -\ii\lambda(e_1^\mu e_2^\nu-e_2^\mu e_1^\nu)
 \Bigr].
 \end{align*}
取空间取向与 $\epsilon^{0123}=+1$ 一致，二维面积形式可协变化为
\begin{equation*}
 \mathsf C^{\mu\nu}(k;n)
 \equiv
 \frac{\epsilon^{\mu\nu\rho\sigma}k_\rho n_\sigma}
 {k\!\cdot n}.
 \end{equation*}
因此每个光子螺旋度的秩一投影算符是
\begin{equation*}
 {
 \mathsf P_\lambda^{\mu\nu}(k;n)
 \equiv
 \varepsilon_\lambda^\mu\varepsilon_\lambda^{\nu *}
 =\frac12\left[
 \mathsf P_\perp^{\mu\nu}(k;n)
 -\ii\lambda\mathsf C^{\mu\nu}(k;n)
 \right].}
 \end{equation*}
立刻有
\begin{equation*}
 \mathsf P_{+}^{\mu\nu}+\mathsf P_{-}^{\mu\nu}
 =\mathsf P_\perp^{\mu\nu}.
 \end{equation*}
因此未偏振光子求和正是两个 Wigner 螺旋度扇区的完备和。
\end{proof}


若振幅写成 $\mathcal M_\lambda=\varepsilon_{\lambda\mu}\mathcal M^\mu$ 且 Ward 条件 $k_\mu\mathcal M^\mu=0$ 已满足，则固定螺旋度的概率核为
\begin{equation}
 |\mathcal M_\lambda|^2
 =\mathcal M_\mu
 \mathsf P_\lambda^{\mu\nu}
 \mathcal M_\nu^*.
 \label{eq:helicity-resolved-photon-amplitude-dp9}
\end{equation}
这里 $n^\mu$ 的改变只改变极化代表元的规范部分；对守恒的振幅，物理结果与 $n^\mu$ 无关。若这一独立性在计算中失败，应先检查 Ward 恒等式，而不能把参考矢量依赖当成可观测偏振效应。

这些投影式已经把常见极化制备统一到同一套密度矩阵语言中：费米子由 \eqnrefs{eq:polarized-spinor-projector-r9,eq:electron-helicity-projector-dp9} 描述，光子由 \eqnrefs{eq:photon-helicity-projector-covariant-dp9} 描述；自旋平均、固定自旋、圆极化或螺旋度不对称只对应不同的实验制备、求和与取差方式，并不需要额外的经验规则。

